\DTMsavedate{mydate}{2022-06-30}
\maketitle{Electrical Measurements}{Electronic Engineering}{\DTMusedate{mydate}}

% ----------------------------------------------------- %
% COURSE INFO
\subsection*{Course Information}
\vspace{0ex}%
\noindent\parbox{0.5\textwidth}{%
    \noindent\begin{tabular}{@{}ll}
                 \textsf{Course:}          & Electrical Measurements \\
                 \textsf{Course Code:}         & 	TEE 1241    \\
                 \textsf{Credit Hours:}    & 	10
    \end{tabular}}
\vspace{2ex}\hrule\vspace{2ex}

% ----------------------------------------------------- %
% INSTRUCTOR INFO
\subsection*{Lecturer Information}
\noindent\parbox{0.5\textwidth}{%
    \noindent\begin{tabular}{@{}ll}
                 \textsf{Name:}         &    Bhekisisa Nyoni         \\
%\textsf{Phone:}        &    \phone          \\
\textsf{Email:}        &    \href{mailto:bhekisisa.dube@nust.ac.zw}{bhekisisa.dube@nust.ac.zw} \\
\textsf{Qualifications:}        &  MPhil in Electronic Engineering (NUST), BEng (Hons) in Electronic Engineering \\
 &                                 (NUST), PGD in Project Management (PMZ), PGD in Higher Education \\
 &                                 (NUST), OSHEMAC (NSSA)

    \end{tabular}}
\vspace{2ex}\hrule\vspace{2ex}


% ----------------------------------------------------- %
% COURSE SYNOPSIS
\subsection*{Course Synopsis}
This course delves into basic measuring devices, ammeters, voltmeters, transducers, accuracy, precision and components.
It looks into electronic measuring instruments, digital voltmeters, multimeters.
It also covers oscilloscopes as measurement instruments, measurement of AC power and measurement of nonelectrical parameters.
\vspace{2ex} \hrule\vspace{2ex}
% ----------------------------------------------------- %
% COURSE TOPICS
\subsection*{Topics}
\begin{outline}
    \1 Fundamentals of Measurement System
    \1 Sensors
    \1 Signal Conditioning
    \1 Analogue Instruments
    \1 Digital Instruments
    \1 Frequency and Waveform Measurement
    \1 Recording Instruments
    \1 Telemetry
    \1 Automatic Measurement Systems
    \1 Errors in Measurement
\end{outline}
\vspace{2ex} \hrule\vspace{2ex}

% ----------------------------------------------------- %
% Students Assenssment
\subsection*{Assessment of Students}
\begin{outline}
    \1 Students are assessed through written assignments and written in-class tests that shall form part of course work assessment mark.
    \1 There shall be a final examination at the end of the semester.
\end{outline}
\vspace{2ex}\hrule\vspace{2ex}

% ----------------------------------------------------- %
% Grade Evaluation
\subsection*{Course Evaluation and Grading Policies}
Grades are based on:
Continuous assessment (\qty{25}{\percent}),
Final Exam (\qty{75}{\percent})
\vspace{2ex}\hrule
\vspace{2ex}\hrule\vspace{2ex}

% ----------------------------------------------------- %
% EQUIPMENT
% https://sites.udel.edu/computing-purchases/personal-specs/
% \subsection*{Essential Equipment}

% \vspace{2ex}\hrule\vspace{2ex}

% Policies and Other information
